(1) 对混合自由能两次微分有
\[
g'=\frac1N\ln\frac{\phi}{1-\phi}+\chi(1-2\phi),\qquad
g''=\frac1{N\phi}+\frac1{N(1-\phi)}-2\chi.
\]
在固定平均成分下小的局部涨落要求 \(g''>0\) 方有局部稳定性，旋节为 \(g''=0\)，故
\[
\boxed{\chi_s(\phi)=\frac1{2N\phi(1-\phi)}}.
\]
进入 \(g''<0\) 区域的均匀混合态对微小浓度波动线性不稳定，而旋节外部仍可能是亚稳而非全球稳定。

(2) 临界点同时要求 \(g''=g'''=0\)。由于
\(g'''=-1/(N\phi^2)+1/[N(1-\phi)^2]\)，在 \(0<\phi<1\) 只有 \(\phi_c=1/2\)，代回得 \(\boxed{\chi_c=2/N}\)。再求 \(g''''=2/(N\phi^3)+2/[N(1-\phi)^3]\)，在临界点为 \(32/N>0\)，因此四阶项正使临界自由能局部有界；与二阶曲率变号组成平均场二阶分相临界结构。

(3) 对称体系 \(g(\phi)=g(1-\phi)\)，若两相共存的A体积分数为 \(\phi_-\) 和 \(\phi_+=1-\phi_-\)，共切线的斜率为零，所以有
\[
\boxed{\frac1N\ln\frac{\phi_-}{1-\phi_-}+\chi(1-2\phi_-)=0}.
\]
令 \(\phi_\pm=1/2\pm m\)，对 \(|m|\ll1\) 展开
\(\ln[(1+2m)/(1-2m)]=4m+16m^3/3+O(m^5)\)，于是
\((4/N-2\chi)m+(16/(3N))m^3+O(m^5)=0\)。非零分支给
\[
\boxed{m^2\simeq\frac{3N}{8}(\chi-2/N)}.
\]
该平方根是对称平均场binodal临界标度，并非旋节方程本身。

(4) 共存线binodal由两相全局自由能共切线决定；旋节线spinodal由单相局部曲率为零决定，两者之间是可能成核的亚稳区。 \(\chi\to0\) 时所有 \(0<\phi<1\) 的混合熵使 \(g''>0\)；对很大 \(N\)，\(\chi_c=2/N\to0\)，长链混合熵变小而更易分相。但无限聚合度、有限链尺度、界面能和非均匀涨落将改变精确转变性质，不能把无梯度均匀自由能曲线冒充真实动力学成核速率。
